\documentclass[10pt]{article}

% ---------- Document Identity (update these only) ----------
\newcommand{\DocStage}{}
\newcommand{\DocTitle}{Your Road to Tier One \textbf{SOF}}
\newcommand{\ProgramName}{182-Week Civilian-to-Selection Readiness Pipeline}
\newcommand{\ProgramSubtitle}{}

% ---------- Page ----------
\usepackage[legalpaper,left=0.5in,right=0.5in,top=0.6in,bottom=0.45in]{geometry}

% ---------- Typography ----------
\usepackage[T1]{fontenc}
\usepackage[utf8]{inputenc}
\usepackage{libertinus}
\usepackage{microtype}
\usepackage{parskip}
\usepackage{ragged2e}
\usepackage{textcomp}
\AtBeginDocument{\color{black}}
\AtBeginDocument{\pagecolor{white}}
\AtBeginDocument{\justifying}

% ---------- Landscape pages ----------
\usepackage{pdflscape}

% ---------- Color ----------
\usepackage[table]{xcolor}
\definecolor{StageBlue}{HTML}{1F4E79}
\definecolor{StageBlueDark}{HTML}{163A5A}
\definecolor{LightGray}{HTML}{F2F4F7}
\definecolor{MidGray}{HTML}{D9DEE5}
\definecolor{RuleGray}{HTML}{C7CED8}
\definecolor{HeaderInk}{HTML}{000000}
\definecolor{Ink}{HTML}{000000}

% ---------- Utility ----------
\usepackage{enumitem}
\sloppy
\setlength{\emergencystretch}{2em}

% ---------- Boxes / UI ----------
\usepackage[most]{tcolorbox}

\tcbset{
  charterTitle/.style={
    enhanced,
    colback=StageBlue!4,
    colframe=StageBlue,
    boxrule=1.25pt,
    arc=3pt,
    left=16pt,right=16pt,top=14pt,bottom=14pt,
    borderline north={3.2pt}{0pt}{StageBlue},
    borderline west={4pt}{0pt}{StageBlue},
    borderline south={1.25pt}{0pt}{StageBlue},
    drop shadow={black!25!white},
  },
  charterRule/.style={
    enhanced,
    colback=LightGray,
    colframe=RuleGray,
    boxrule=0.85pt,
    arc=3pt,
    left=12pt,right=12pt,top=10pt,bottom=10pt,
    borderline west={3pt}{0pt}{StageBlue},
  },
  charterBadge/.style={
    enhanced,
    colback=StageBlue,
    coltext=white,
    colframe=StageBlueDark,
    boxrule=0.6pt,
    arc=2pt,
    left=6pt,right=6pt,top=3pt,bottom=3pt,
  }
}
\newtcbox{\Badge}{nobeforeafter,charterBadge}

% ---------- Lists ----------
\setlist[itemize]{leftmargin=1.5em, itemsep=0.25em, topsep=0.25em}
\setlist[enumerate]{leftmargin=1.8em, itemsep=0.25em, topsep=0.25em}

% ---------- TOC ----------
\usepackage{tocloft}
\setcounter{tocdepth}{3}
\setlength{\cftbeforesecskip}{3pt}
\setlength{\cftbeforesubsecskip}{2pt}
\setlength{\cftbeforesubsubsecskip}{0pt}
\renewcommand{\cftsecfont}{\bfseries}
\renewcommand{\cftsecpagefont}{\bfseries}
\renewcommand{\cftsubsecfont}{\normalfont}
\renewcommand{\cftsubsecpagefont}{\normalfont}
\renewcommand{\cftsubsubsecfont}{\normalfont}
\renewcommand{\cftsubsubsecpagefont}{\normalfont}
\setlength{\cftsecindent}{0pt}
\setlength{\cftsubsecindent}{1.25em}
\setlength{\cftsubsubsecindent}{2.8em}
\setlength{\cftsecnumwidth}{2.1em}
\setlength{\cftsubsecnumwidth}{2.8em}
\setlength{\cftsubsubsecnumwidth}{3.6em}

% Remove the built-in TOC heading so only the custom heading is shown
\makeatletter
\renewcommand{\tableofcontents}{\@starttoc{toc}}
\makeatother

% ---------- Header/Footer: page number only ----------
\usepackage{fancyhdr}
\pagestyle{fancy}
\fancyhf{}
\renewcommand{\headrulewidth}{0pt}
\renewcommand{\footrulewidth}{0pt}
\fancyfoot[C]{\thepage}

% ---------- Section formatting ----------
\usepackage{titlesec}

\titleformat{\section}
  {\Large\bfseries\color{Ink}}
  {\thesection}{0.6em}{}
  [\vspace{0.25em}\textcolor{StageBlue}{\rule{\linewidth}{0.8pt}}]

\titleformat{\subsection}
  {\large\bfseries\color{Ink}}
  {\thesubsection}{0.6em}{}

\titleformat{\subsubsection}
  {\normalsize\bfseries\color{Ink}}
  {\thesubsubsection}{0.6em}{}

\titlespacing*{\section}{0pt}{0.95em}{0.35em}
\titlespacing*{\subsection}{0pt}{0.65em}{0.25em}
\titlespacing*{\subsubsection}{0pt}{0.55em}{0.2em}

% ---------- Table engine ----------
\usepackage{tabularray}
\UseTblrLibrary{booktabs}

% ---------- tabularray: GLOBAL table treatment ----------
\SetTblrInner{
  cells = {fg=black, bg=white, valign=t},
}

% ---------- Header helpers ----------
\newcommand{\Hdr}[1]{\shortstack[c]{\strut #1\strut}}

% ---------- Lines ----------
\newcommand{\TitleLine}{\textcolor{StageBlue}{\rule{0.98\linewidth}{0.95pt}}}
\newcommand{\SubTitleLine}{\textcolor{RuleGray}{\rule{0.98\linewidth}{0.65pt}}}

% ---------- Continued on next page ----------
\newcommand{\ContinuedOnNextPageInline}{%
  \par
  \vfill
  \noindent\textcolor{RuleGray}{\rule{\linewidth}{1pt}}\vspace{-2pt}
  \noindent\hfill\textit{Continued on next page}\par\vspace{-10pt}
  \noindent\textcolor{RuleGray}{\rule{\linewidth}{1pt}}
  \clearpage
}

% ---------- PDF hyperlinks ----------
\usepackage[hidelinks]{hyperref}
\hypersetup{
  colorlinks=false,
  pdfborder={0 0 0},
  linktoc=all,
  pdftitle={\DocTitle},
  pdfauthor={\ProgramName}
}

% =========================================================
\begin{document}

\section*{Stage 11.B — Final Preparation and Test Optimization}
\phantomsection
\addcontentsline{toc}{section}{Stage 11.B — Final Preparation and Test Optimization}

\subsection*{11.B.1 Purpose and Scope}
\phantomsection
\addcontentsline{toc}{subsection}{11.B.1 Purpose and Scope}

11.B defines the final preparation and test optimization posture for the selection approach period. In this system, “optimization” means maximizing validity and repeatability of gate-relevant tests during the final phases by enforcing standardization, controlling variables, preventing invalid tests, and reducing unsafe attempts. It exists to ensure that gate outcomes (\textbf{ADVANCE} / \textbf{HOLD} / \textbf{REGRESS} / \textbf{MEDICAL} \textbf{HOLD}) are based on admissible \textbf{VALID} evidence only and that failures are handled through reslot discipline rather than improvisation.

Exclusions:

\begin{itemize}
\item No programming detail (no workouts, sets/reps, intervals, weekly splits, day-by-day schedules).
\item No new tests, protocols, thresholds, domains, phase purposes, or equipment categories.
\item No medical diagnosis/treatment guidance; no prescription or hormone guidance; no buyer’s guide content.
\end{itemize}

\newpage
\subsection*{11.B.2 Test Optimization Doctrine — Validity First}
\phantomsection
\addcontentsline{toc}{subsection}{11.B.2 Test Optimization Doctrine — Validity First}

\begin{itemize}
\item “Optimization” \textbf{MUST} be defined as standardization, repeatability, and validity protection, not hype, risk escalation, or last-minute intensity chasing.
\item Gate evidence \textbf{MUST} be admissible:
\begin{itemize}
\item Only \textbf{VALID} tests provide gate credit.
\item \textbf{INVALID} tests provide no gate credit and are treated as not yet tested for gate decisions.
\end{itemize}
\item Replication posture \textbf{MUST} follow the established upstream rule where it is used; do not invent new replication rules.
\item Pre-test readiness screening \textbf{MUST} be pass/fail and \textbf{MUST} trigger \textbf{STOP} \textbf{NOW} and/or \textbf{MEDICAL} \textbf{HOLD} posture when indicated.
\item No novelty near gates/test weeks is mandatory:
\begin{itemize}
\item no new footwear system, no new route class, no new tool system, no new caffeine pattern, no new supplement category, no new electrolyte practice inside the protected novelty windows.
\end{itemize}
\item When a variable change is unavoidable due to logistics:
\begin{itemize}
\item it \textbf{MUST} be logged as a controlled-variable deviation,
\item outcomes \textbf{MUST} be interpreted conservatively and may require reslot to preserve admissibility.
\end{itemize}
\end{itemize}

\begingroup
\scriptsize
\setlength{\tabcolsep}{2pt}
\renewcommand{\arraystretch}{1.12}

\begin{longtblr}{
width=\linewidth,
colspec = {
X[l,1.35]
X[l,1.55]
X[l,2.05]
X[l,2.05]
X[l,2.20]
},
rowhead=1,
hlines,
vlines,
cells = {hsep=2pt, vsep=6pt, valign=m},
row{1} = {bg=StageBlue!15, fg=black, font=\bfseries, valign=m},
row{odd} = {bg=white},
row{even} = {bg=LightGray},
}
\Hdr{Optimization Lever} &
\Hdr{What It Means} &
\Hdr{Allowed} &
\Hdr{Not Allowed} &
\Hdr{Validity Risk if Violated} \\
Environment control &
Same Route/Pool/Machine identity and conditions posture &
Same Environment Unit \textbf{ID}; consistent surface class; documented conditions &
Unmeasured route drift; switching pools/lanes without logging &
Results become non-comparable or inadmissible due to missing IDs/conditions \\
Tool standardization &
Stable Timer/Scale/Tape and stable methods &
Same tool identity; backups ready; settings logged &
Tool swaps in protected windows without logging &
Tool drift contaminates comparability; can render evidence inadmissible \\
Warm-up completion &
Required, logged, consistent &
Standardized warm-up completed and recorded &
Skipping or improvising warm-up on test day &
Increased injury risk; protocol non-compliance can invalidate test \\
Controlled variables &
Stable inputs near gates &
Stable foods/hydration/caffeine routine; stable sleep routine &
Novel foods, new caffeine pattern, new supplement category, new electrolyte routine &
Confounds performance; increases GI and readiness variance; validity contamination \\
Spacing/stacking discipline &
Stage 3 anti-stacking intent preserved &
Staged battery posture; reslot rather than compress &
Make-up stacking to “hit deadlines” &
Injury risk increases; fatigue contamination invalidates evidence \\
Stop posture execution &
Safety overrides schedule &
\textbf{STOP} \textbf{NOW} executed immediately when triggered &
“Push through” symptoms to complete test &
Unsafe exposure; evidence inadmissible for advancement \
\end{longtblr}

\endgroup

\newpage

\subsection*{11.B.3 Pre-Test Standardization Checklist — Universal}
\phantomsection
\addcontentsline{toc}{subsection}{11.B.3 Pre-Test Standardization Checklist — Universal}

Universal test-day checklist (applies to all tests in all domains). Pass/Fail is operational; failures trigger reslot or substitution per Stage 3 and Stage 2.E posture.

\begingroup
\scriptsize
\setlength{\tabcolsep}{2pt}
\renewcommand{\arraystretch}{1.12}

\begin{longtblr}{
width=\linewidth,
colspec = {
X[l,3.05]
Q[c,wd=0.85in]
X[l,1.85]
X[l,2.85]
},
rowhead=1,
hlines,
vlines,
cells = {hsep=2pt, vsep=6pt, valign=m},
row{1} = {bg=StageBlue!15, fg=black, font=\bfseries, valign=m},
row{odd} = {bg=white},
row{even} = {bg=LightGray},
}
\Hdr{Item} &
\Hdr{Pass/Fail} &
\Hdr{Evidence pointer or log field} &
\Hdr{If Fail Action} \\
Environment Unit \textbf{ID} selected and confirmed (\textbf{ROUTE} / \textbf{POOL} / \textbf{MACH}) &
&
Stage 2.E Test Event Log: Environment Unit \textbf{ID} &
Do not test; reslot to next protected window; if training exposure occurs, log as non-gate attempt with full context \\
Tool standard confirmed (Timer / Scale / Tape; tool identity recorded; backup ready where applicable) &
&
Stage 2.E Test Event Log: Tool Standard; tool description/\textbf{ID} &
Do not test gate-grade; restore tool readiness; reslot \\
Warm-up completed and logged &
&
Stage 2.E: \textbf{WARMUP} field &
If not completed: complete warm-up before test; if cannot: abort attempt and reslot \\
Evidence artifact posture satisfied where required by protocol &
&
Stage 2.E: Evidence Ref field; protocol artifact requirement field &
If missing: tag test \textbf{INVALID} for gate use; reslot; do not claim credit \\
No novelty near gates confirmed (food/caffeine/supplements/footwear/tools) &
&
Stage 2.E Notes + nutrition log: novelty Y/N &
If novelty present inside window: treat as confounder; default reslot for gate-grade attempt; log deviation \\
Sleep and recovery readiness checked (trend-based, non-medical) &
&
Stage 2 daily/session logs: sleep hrs/quality; fatigue/soreness &
If degraded trend: convert to check-in or deload-only; reslot maximal tests \\
Stop-rule screen passed (no red flags; no gait-altering pain) &
&
Stage 1.F posture; Stage 2.E \textbf{SAFETY} field; pain/gait flags &
If failed: \textbf{STOP} \textbf{NOW}; \textbf{MEDICAL} \textbf{HOLD} if indicated; reslot all tests \
\end{longtblr}

\endgroup

\newpage
\subsection*{11.B.4 Domain-Specific Test-Week Optimization Rules}
\phantomsection
\addcontentsline{toc}{subsection}{11.B.4 Domain-Specific Test-Week Optimization Rules}

Domain posture is operational and conservative; it prevents invalidators, enforces logging completeness, and preserves Stage 3 anti-stacking intent without creating day-by-day schedules. Stage 2 test IDs are referenced directly for core tests in this package.

Domain rules (high-level):

\begin{itemize}
\item \textbf{CAL}:
\begin{itemize}
\item Control variables: surface, camera/artifact posture where required, counting authority posture.
\item Prevent invalidators: inconsistent standards, missing artifacts, incomplete warm-up.
\end{itemize}
\item \textbf{RUN}:
\begin{itemize}
\item Control variables: Route \textbf{ID} or Machine \textbf{ID} + settings; footwear stability; weather/surface class logging.
\item Prevent invalidators: unsafe conditions, unlogged treadmill settings, route drift.
\end{itemize}
\item \textbf{RUCK}:
\begin{itemize}
\item Control variables: Route \textbf{ID}; dry load verification; carried water logged separately; ruck fit stability; footwear + sock system stability.
\item Prevent invalidators: load shift, blister progression, unsafe weather/visibility.
\end{itemize}
\item \textbf{SWIM}:
\begin{itemize}
\item Control variables: Pool \textbf{ID} + pool length unit; lane conditions; no fins; stroke rules per protocol; rest governance.
\item Prevent invalidators: missing Pool \textbf{ID}/length unit, standing/hanging violations, unsafe pool conditions, illness.
\end{itemize}
\item \textbf{STR}:
\begin{itemize}
\item Control variables: implement identity; load method; submax proxy posture; safety discipline.
\item Prevent invalidators: unsafe attempts, missing tool identity/load method, novelty implement changes in protected windows.
\end{itemize}
\item \textbf{BODYCOMP} / \textbf{RECOVERY} / \textbf{DURABILITY} monitoring:
\begin{itemize}
\item Primary function: validity protection and readiness constraint signals; not a substitute for performance tests.
\item Prevent invalidators: missing logs, backfilled entries, method drift for BF% and circumference landmarks.
\end{itemize}
\end{itemize}

\begingroup
\scriptsize
\setlength{\tabcolsep}{2pt}
\renewcommand{\arraystretch}{1.12}

\begin{longtblr}{
width=\linewidth,
colspec = {
Q[l,wd=1.20in]
X[l,2.35]
X[l,2.35]
X[l,2.05]
X[l,2.05]
},
rowhead=1,
hlines,
vlines,
cells = {hsep=2pt, vsep=6pt, valign=m},
row{1} = {bg=StageBlue!15, fg=black, font=\bfseries, valign=m},
row{odd} = {bg=white},
row{even} = {bg=LightGray},
}
\Hdr{Domain} &
\Hdr{Common Invalidators} &
\Hdr{Prevention Controls} &
\Hdr{Reslot Trigger} &
\Hdr{Evidence to Log} \\
\textbf{CAL} &
Missing required artifact/evidence ref; inconsistent form standard; incomplete warm-up &
Confirm artifact capture and Evidence Ref before start; use stable surface and standard; warm-up logged &
Artifact missing or non-compliant; form standard violated; stop-rule event &
Test Validity Tag; Evidence Ref; Tool Standard (timer); conditions notes \\
\textbf{RUN} &
Missing Route/Machine \textbf{ID}; treadmill settings not logged; unsafe ice/heat/air quality; route drift &
Register Route \textbf{ID}; log surface class; if treadmill, log Machine \textbf{ID} + speed/grade; weather check &
Unsafe conditions; missing \textbf{ID}/settings; gait-altering pain &
Route/Machine \textbf{ID}; Tool Standard; conditions; stop/abort; footwear note \\
\textbf{RUCK} &
Dry load not verified; water not separated; load shift; blister progression; unsafe visibility &
Dry load recorded; water logged separately; verify fit; visibility kit when needed &
Load verification missing; gait-altering pain; hotspot→blister trend; unsafe conditions &
Route \textbf{ID}; dry load; water carried; conditions; foot flags; stop/abort \\
\textbf{SWIM} &
Missing Pool \textbf{ID}/length unit; rule violations; pool crowding/closure; illness symptoms &
Confirm Pool \textbf{ID} + length unit; lane conditions notes; stable routine &
Pool \textbf{ID}/unit missing; rule violation; unsafe pool condition; illness &
Pool \textbf{ID}; length unit; conditions; Evidence Ref; validity tag \\
\textbf{STR} &
Implement identity drift; unsafe attempt; missing load method &
Log implement type and load method; maintain submax proxy posture &
Unsafe execution; missing required fields &
Tool identity; load method; validity tag; notes \\
\textbf{BODYCOMP} / \newline \textbf{RECOVERY} /\newline \textbf{DURABILITY} &
Missing weekly averages; missing pain/foot flags; method drift &
Enforce same-day capture; mark \textbf{MISSING} rather than infer &
Missing required fields before gate decisions &
Weekly \textbf{BW} avg; circumference; sleep/fatigue/soreness; pain/gait/foot flags \
\end{longtblr}

\endgroup

\subsection*{11.B.5 Deload and Test Window Execution Posture}
\phantomsection
\addcontentsline{toc}{subsection}{11.B.5 Deload and Test Window Execution Posture}

\begin{itemize}
\item Deload/Test week equals test week posture per Stage 3:
\begin{itemize}
\item primary goal is admissible evidence collection under reduced fatigue and controlled variables.
\end{itemize}
\item Avoid stacking beyond Stage 3 principles:
\begin{itemize}
\item staged batteries are permitted, but they must preserve spacing intent and not compress high-risk domains.
\end{itemize}
\item Preserve safe cadence posture:
\begin{itemize}
\item maintain the 6-session rhythm where feasible while downshifting density and risk; do not remove required session elements.
\end{itemize}
\item Staging and logistics posture (non-programming):
\begin{itemize}
\item confirm equipment readiness before the protected week (no last-minute changes),
\item confirm route/pool access and environment feasibility,
\item confirm tool readiness and backup posture,
\item confirm hydration and nutrition stability (Stage 10) with no novelty.
\end{itemize}
\item Travel contingency posture:
\begin{itemize}
\item if travel threatens feasibility, reslot tests rather than compress; do not stack make-ups.
\end{itemize}
\end{itemize}

Staged battery guidance (no day-by-day scheduling):

\begin{itemize}
\item A staged battery is a protected-week intent that sequences major domains without implying same-day stacking.
\item The intent \textbf{MUST} preserve:
\begin{itemize}
\item spacing between maximal \textbf{RUN} and maximal \textbf{RUCK},
\item controlled pool scheduling for \textbf{SWIM} separate from other high-cost exposures when feasible,
\item reduced fatigue and stable inputs for admissibility.
\end{itemize}
\end{itemize}

\newpage
\subsection*{11.B.6 Reslot and Invalidation Handling — Stage 3 Alignment}
\phantomsection
\addcontentsline{toc}{subsection}{11.B.6 Reslot and Invalidation Handling — Stage 3 Alignment}

\begin{itemize}
\item \textbf{INVALID} test equals no gate credit:
\begin{itemize}
\item treat as not yet tested for gate decisions.
\end{itemize}
\item Reslot to next protected window per Stage 3:
\begin{itemize}
\item reslot is the default; do not compress or stack missed tests.
\end{itemize}
\item No stacking make-ups:
\begin{itemize}
\item missed tests are not recovered by packing multiple high-cost tests into one week.
\end{itemize}
\item Default gate outcome posture when evidence is missing:
\begin{itemize}
\item \textbf{HOLD} is the default when required evidence is missing or inadmissible.
\item \textbf{MEDICAL} \textbf{HOLD} is used only when explicit triggers apply per Stage 1.F (illness/red flags/clinician restrictions).
\end{itemize}
\item Environment invalidation examples that require reslot or indoor substitution (only if comparability preserved):
\begin{itemize}
\item unsafe footing/ice glaze,
\item lightning risk,
\item poor air quality warnings,
\item pool closure or crowding that breaks validity/safety,
\item extreme heat/cold triggers per Stage 10.C and environment governance posture.
\end{itemize}
\item Known high-consequence timing constraints that affect reslot feasibility:
\begin{itemize}
\item underwater governance availability and admissibility controls (\textbf{ISS-08-002}; \textbf{ISS-08-003}),
\item long ruck compounded dependency stack (\textbf{ISS-08-004}).
\end{itemize}
\end{itemize}

Requires Stage 8 review flags (if encountered):

\begin{itemize}
\item Any phase card that schedules a decisive gate test outside protected windows.
\item Any gate reliance on a standard whose enabling capability is not available per Stage 6 timing (e.g., pull-up feasibility timing conflict: \textbf{ISS-08-001}).
\end{itemize}

\newpage
\subsection*{11.B.7 Capability Readiness and Tooling Readiness}
\phantomsection
\addcontentsline{toc}{subsection}{11.B.7 Capability Readiness and Tooling Readiness}

\begin{itemize}
\item Required capabilities \textbf{MUST} be present by the phase timing established in Stage 6:
\begin{itemize}
\item do not attempt a gate test without required Tier 1 resources when applicable.
\end{itemize}
\item Do not “make do” with unsafe substitutes:
\begin{itemize}
\item if capability is missing, default to \textbf{HOLD} and reslot; preserve safety and admissibility.
\end{itemize}
\item Substitution posture:
\begin{itemize}
\item substitute only when it preserves validity and is permitted by protocol posture,
\item document substitutions per Stage 2.E with environment IDs, tool standard, conditions notes, and reason posture.
\end{itemize}
\item Known readiness blockers (must be handled explicitly, not silently):
\begin{itemize}
\item pull-up feasibility timing risk (\textbf{ISS-08-001}),
\item underwater governance feasibility/admissibility (\textbf{ISS-08-002}; \textbf{ISS-08-003}),
\item long ruck dependency stack (\textbf{ISS-08-004}).
\end{itemize}
\end{itemize}

\begingroup
\scriptsize
\setlength{\tabcolsep}{2pt}
\renewcommand{\arraystretch}{1.12}

\begin{longtblr}{
width=\linewidth,
colspec = {
Q[l,wd=1.75in]
X[l,2.55]
X[l,2.10]
X[l,2.60]
},
rowhead=1,
hlines,
vlines,
cells = {hsep=2pt, vsep=6pt, valign=m},
row{1} = {bg=StageBlue!15, fg=black, font=\bfseries, valign=m},
row{odd} = {bg=white},
row{even} = {bg=LightGray},
}
\Hdr{Capability Area} &
\Hdr{Readiness Check} &
\Hdr{Evidence Pointer} &
\Hdr{If Fail Action} \\
Timing/measurement tools &
Timer/scale/tape and route/pool/machine \textbf{ID} readiness verified &
Stage 2.E logs; 7.18 overview; binder pointer to tool registry &
\textbf{HOLD} gate attempt; restore tool readiness; reslot \\
Route readiness &
Route \textbf{ID} registered; surface class stable; safety feasible &
Stage 2.E Route \textbf{ID}; conditions notes &
Substitute indoors only if comparable; otherwise reslot \\
Pool readiness &
Pool \textbf{ID} + length unit stable; access confirmed &
Stage 2.E Pool \textbf{ID}; 7.17 overview &
Reslot swim test; do not infer distance/unit \\
Foot system readiness &
Footwear + sock system + hotspot controls stable &
7.11; foot flags in logs &
Downshift exposure; \textbf{HOLD}/risk posture; reslot high-cost tests \\
Ruck system readiness &
Ruck fit and load stability verified; dry load verifiable &
7.07; Stage 2.E load logs &
Do not ruck test; \textbf{HOLD} and reslot; correct readiness gap \\
\textbf{CAL} artifact capture &
Evidence ref posture ready where required &
Stage 2.E Evidence Ref; capture plan in 7.16/7.18 &
If missing, test is \textbf{INVALID} for gate use; reslot \\
Underwater governance readiness (if applicable) &
Lifeguard + dedicated spotter + controlled pool confirmed; otherwise prohibited &
Governance confirmation logged; 7.17 posture &
Not scheduled; inadmissible without governance; reslot; \textbf{ISS-08-002}/003 handling \
\end{longtblr}


\endgroup

Requires Stage 8 review flags (if encountered):

\begin{itemize}
\item Any phase that demands a capability earlier than Stage 6 introduces it at Tier 1+.
\item Any underwater exposure implied as routine rather than conditional and supervised-only.
\end{itemize}

\newpage
\subsection*{11.B.8 Compliance Checklist — Pass/Fail}
\phantomsection
\addcontentsline{toc}{subsection}{11.B.8 Compliance Checklist — Pass/Fail}

\begin{itemize}
\item \textbf{PASS}/\textbf{FAIL}: No programming detail included (no workouts, sets/reps, intervals, weekly splits, day-by-day schedules).
\item \textbf{PASS}/\textbf{FAIL}: No new tests, thresholds, domains, protocols, or phase purposes introduced.
\item \textbf{PASS}/\textbf{FAIL}: Validity tokens and gate outcomes use Stage 1.F terms exactly: \textbf{VALID} / \textbf{MODIFIED} / \textbf{INVALID}; \textbf{VALID} / \textbf{INVALID}; \textbf{ADVANCE} / \textbf{HOLD} / \textbf{REGRESS} / \textbf{MEDICAL} \textbf{HOLD}; \textbf{STOP} \textbf{NOW} / \textbf{MEDICAL} \textbf{HOLD}.
\item \textbf{PASS}/\textbf{FAIL}: Environment Unit \textbf{ID} and Tool Standard posture enforced; Stage 2.E required fields treated as mandatory; missing fields recorded as \textbf{UNKNOWN}/\textbf{MISSING} with no inference and treated as inadmissible when required.
\item \textbf{PASS}/\textbf{FAIL}: No novelty near gates posture stated and enforced (including caffeine pattern change restriction window).
\item \textbf{PASS}/\textbf{FAIL}: Reslot posture follows Stage 3 and no stacking is explicit; invalid tests provide no gate credit.
\item \textbf{PASS}/\textbf{FAIL}: Capability readiness posture aligns to Stage 6 timing; no gate attempt proceeds without required Tier 1 readiness where applicable.
\item \textbf{PASS}/\textbf{FAIL}: Any detected mismatch is flagged "Requires Stage 8 review" in 11.B.6 and 11.B.7 rather than silently fixed.
\end{itemize}

\end{document}